\documentclass[twocolumn]{article}
\usepackage{amsmath}
\usepackage{booktabs}
\usepackage{longtable}
\usepackage[ruled,linesnumbered]{algorithm2e}
\usepackage{subcaption}
\usepackage{graphicx}
\title{Tables, Algorithms and Subfigures}
\author{Test Author}
\begin{document}
\maketitle
\section{Introduction}
We compare several baselines~\cite{smith2020}.
\begin{table}[t]
\centering
\caption{Success rates on three tasks.}
\begin{tabular}{lccc}
\toprule
Method & Walk & Run & Jump \\
\midrule
Ours & 98.1 & 95.4 & 91.2 \\
Baseline & 92.3 & 88.0 & 80.5 \\
\bottomrule
\end{tabular}
\end{table}
\begin{algorithm}[t]
\caption{Training procedure}
\KwIn{Learning rate $\alpha$, episodes $N$}
\KwOut{Policy $\pi_\theta$}
\For{$i \leftarrow 1$ \KwTo $N$}{
  Collect rollouts and update $\theta$ by gradient ascent\;
}
\Return $\pi_\theta$\;
\end{algorithm}
\begin{figure}[t]
\centering
\begin{subfigure}{0.45\linewidth}\centering
\rule{2cm}{1cm}
\caption{Simulation.}
\end{subfigure}
\begin{subfigure}{0.45\linewidth}\centering
\rule{2cm}{1cm}
\caption{Hardware.}
\end{subfigure}
\caption{Qualitative comparison of the learned gait.}
\end{figure}
\bibliographystyle{plain}
\bibliography{refs}
\end{document}
